\section*{Bab \ref{ch:b2:climate-variability} --- Variabilitas Iklim}

\begin{solution}{exo:b2:climate-variability:1}
(a) $29.27 - 26.71 = \qty{2.56}{\celsius}$ pada 2026; $28.79 - 26.73 = \qty{2.06}{\celsius}$ pada 1997.
(b) September 2026 adalah bulan yang lebih hangat, sebesar \qty{0.48}{\celsius}, sekaligus anomali yang lebih besar, sebesar
\qty{0.50}{\celsius}. (c) Setiap tahun dibandingkan dengan periode dasar 30 tahun yang berpusat padanya, yang digeser
setiap lima tahun; \omterm{def:b2:climate-variability:modes}{klimatologi} September tahun 1980-an dan 1990-an sedikit berbeda dari \omterm{def:b2:climate-variability:modes}{klimatologi}
tahun 2010-an dan 2020-an.
\end{solution}

\begin{solution}{exo:b2:climate-variability:2}
(a) JJA: $(0.79 + 0.95 + 1.28)/3 = 1.007$; JAS: $(0.95 + 1.28 + 1.55)/3 = 1.260$; ASO: $(1.28 +
1.55 + 1.68)/3 = 1.503$; SON: $(1.55 + 1.68 + 2.01)/3 = 1.747$; OND: $(1.68 + 2.01 + 2.02)/3 =
1.903$, semuanya dalam \unit{\celsius}. (b) Dibulatkan ke satu desimal: 1.0, 1.3, 1.5, 1.7, 1.9, yaitu nilai-nilai
dalam tabel. (c) Pada ASO 2023, dengan 1.5.
\end{solution}

\begin{solution}{exo:b2:climate-variability:3}
(a) Indeks itu mengikuti simpangan tekanan di Tahiti dikurangi tekanan di Darwin: di bawah rata-rata
di Tahiti dan di atas rata-rata di Darwin, indeks itu negatif. (b) El~Niño. (c) Angin pasat timuran yang lebih lemah, bahkan
angin baratan dekat permukaan di barat, dan hujan yang lebih sedikit daripada biasanya di atas Indonesia, karena
konveksi dalam telah bergeser ke timur.
\end{solution}

\begin{solution}{exo:b2:climate-variability:4}
(a) $804.1/880.6 = \qty{91.3}{\percent}$ dari rata-rata: di dalam rentang 90 sampai \qty{110}{\percent}, sebuah
musim normal, meskipun di bagian bawahnya. (b) Total musim itu $556 + 659 + 427 + 252 =
\qty{1894}{mm}$; Juli membawa $659/1894 = \qty{34.8}{\percent}$ darinya, lebih dari sepertiga.
\end{solution}

\begin{solution}{exo:b2:climate-variability:5}
(a) $\tau = 1025 \times 3992 \times 50/16.1 = \qty{1.271e7}{s} = \qty{147.1}{\day}$ untuk
\qty{50}{m}, dan dua kalinya, \qty{294.2}{\day}, untuk \qty{100}{m}. (b) $\mathrm{e}^{-30.44/147.1} =
0.813$ dan $\mathrm{e}^{-30.44/294.2} = 0.902$. (c) Waktu redam sebanding dengan $h$: lapisan
campuran yang lebih dangkal menyimpan lebih sedikit kalor untuk fluks yang sama dan melupakan anomali lebih cepat, sehingga
ingatan samudra, dan keterprakiraan yang diberikannya, menurun.
\end{solution}

\begin{solution}{exo:b2:climate-variability:6}
(a) Tiga bulan adalah 91.3 hari: $1.2 \times \mathrm{e}^{-91.3/147.1} = \qty{0.65}{K}$; enam bulan
adalah 182.6 hari: $1.2 \times \mathrm{e}^{-182.6/147.1} = \qty{0.35}{K}$. (b) Anomali itu meluruh
dengan lambat, sehingga menjadi tebakan anomali bulan depan yang lebih baik daripada nol. Tetapi gaya paksa acak menambahkan
anomali baru setiap hari, yang tidak diketahui pada saat prakiraan dibuat; seiring meluruhnya anomali awal,
anomali-anomali baru itu mengambil bagian yang makin besar, dan prakiraan itu bergeser menuju \omterm{def:b2:climate-variability:modes}{klimatologi}, tempat ia tidak
lagi membantu.
\end{solution}

\begin{solution}{exo:b2:climate-variability:7}
(a) $160 \times 86\,400 = \qty{1.38e7}{J}$, \qty{13.8}{MJ} per meter persegi dan per hari. (b) Lapisan
itu memuat $(1000 - 850) \times 100/9.81 = \qty{1529}{kg}$ udara per meter persegi, dengan kapasitas
kalor $1529 \times 1004.5 = \qty{1.536e6}{J.m^{-2}.K^{-1}}$; lapisan itu akan menghangat $1.38 \times
10^7/1.536 \times 10^6 = \qty{9.0}{K}$ sehari. (c) Lapisan itu meradiasikan kalor dalam inframerah, ke
antariksa dan ke udara di atasnya; dan udara hangat naik serta bercampur dengan lapisan di atasnya, atau
dibawa pergi oleh angin, termasuk aliran masuk udara yang lebih sejuk yang ditarik oleh depresi termal itu sendiri.
\end{solution}

\begin{solution}{exo:b2:climate-variability:8}
(a) $(-1.93 - 1.11 - 1.98)/3 = -1.673$, sekitar $-1.67$. (b) Itulah musim dingin terendah dalam catatan.
(c) NAO yang sangat negatif: Depresi Islandia dan Antisiklon Azores yang lemah, arus jet yang lebih lemah, dan jalur badai
yang lebih ke selatan daripada biasanya; musim dingin yang dingin dan kering di Eropa utara, musim dingin yang basah di Eropa selatan,
dan lebih banyak serbuan udara dingin di Amerika Serikat bagian timur.
\end{solution}

\begin{solution}{exo:b2:climate-variability:9}
(a) $9.3 - 0.4 = \qty{8.9}{mm}$ sehari, lebih dari 2.5: termasuk wilayah itu. (b) Di Belahan Bumi Selatan
musim panas jatuh pada Desember--Februari: $6.1 - 4.2 = \qty{1.9}{mm}$ sehari, kurang dari 2.5: di luar wilayah itu,
meskipun hujan di sana lebih banyak daripada di kebanyakan tempat dalam wilayah tersebut. (c) $1894/122 = \qty{15.5}{mm}$
sehari; Kerala baru akan berada di luar wilayah itu jika laju musim dinginnya melebihi $15.5 - 2.5 = \qty{13.0}{mm}$ sehari.
(d) \omterm{def:b2:climate-variability:monsoon}{Monsun} adalah suatu pembalikan musiman: tempat basah yang basah sepanjang tahun, seperti stasiun
kedua, tidak memiliki \omterm{def:b2:climate-variability:monsoon}{monsun}, sedangkan tempat kering dengan musim hujan yang singkat dan intens memilikinya.
\end{solution}

\begin{solution}{exo:b2:climate-variability:10}
(a) Dengan menurunkan, $\ddot T = a\dot H = -abT$: persamaan osilator harmonik dengan frekuensi
sudut $\omega = \sqrt{ab}$, dan periode $2\pi/\sqrt{ab}$. (b) $2\pi/\sqrt{ab} = 4$ memberikan $ab =
(2\pi/4)^2 = \qty{2.47}{yr^{-2}}$. (c) Dengan $T = u\,\mathrm{e}^{-rt}$, $H = v\,\mathrm{e}^{-rt}$,
$\dot T = (\dot u - ru)\,\mathrm{e}^{-rt} = (av - ru)\,\mathrm{e}^{-rt}$, sehingga $\dot u = av$, dan
demikian pula $\dot v = -bu$: $(u, v)$ berosilasi tanpa redaman, dan $(T, H)$ adalah osilasi itu
dikalikan dengan $\mathrm{e}^{-rt}$. Dalam satu periode amplitudo turun dengan faktor $\mathrm{e}^{-0.25 \times
4} = \mathrm{e}^{-1} = 0.37$. (d) Suatu gaya paksa: fluktuasi acak angin pasat, yang
mengentak sistem itu berulang kali, dan \omterm{def:b2:climate-variability:bjerknes}{umpan balik Bjerknes}, yang membuat peredaman lebih lemah atau
bahkan negatif.
\end{solution}

\begin{solution}{exo:b2:climate-variability:11}
(a) $A = R^2 \,\Delta\lambda\,(\sin\ang{5} - \sin(\ang{-5})) = (6.371 \times 10^6)^2 \times 0.8727
\times 0.1743 = \qty{6.17e12}{m^2}$. (b) $Q = 1025 \times 3992 \times 100 \times 2 \times 6.17
\times 10^{12} = \qty{5.05e21}{J}$. (c) Enam bulan adalah $182.6 \times 86\,400 = \qty{1.578e7}{s}$;
per meter persegi, $8.18 \times 10^8/1.578 \times 10^7 = \qty{51.9}{W/m^2}$, yang dipertahankan selama setengah
tahun. (d) Dari samudra itu sendiri: angin pasat yang lebih lemah mengangkat lebih sedikit air dingin di timur, dan air
hangat dari barat, yang tersimpan di atas termoklin yang lebih dalam, menyebar ke timur; sebagian besar penghangatan itu
dapat berasal dari kalor yang sudah ada di samudra, yang dipindahkan dan bukan diperoleh.
\end{solution}

\begin{solution}{exo:b2:climate-variability:12}
(a) Lingkar itu mengembalikan $0.3g$, lalu $0.3g^2$, dan seterusnya: totalnya $0.3(1 + g + g^2 + \dots) =
0.3/(1 - g)$, deret geometri yang konvergen untuk $g < 1$: \qty{0.6}{K}, \qty{1.5}{K}, dan
\qty{3.0}{K} untuk $g = 0.5$, 0.8, dan 0.9. (b) Deret itu divergen: anomali akan tumbuh tanpa
batas; pertumbuhan itu harus dihentikan oleh proses-proses yang tidak dimasukkan dalam lingkar, yaitu peredaman oleh fluks kalor
permukaan dan pengurasan simpanan air hangat yang dijelaskan oleh osilator pengisian ulang. (c)
Makin hangat permukaan laut, makin mudah penghangatan kecil memicu konveksi dalam dan hujan,
yang pemanasannya menggerakkan angin; itulah musim ketika dampak ENSO paling kuat.
\end{solution}

\begin{solution}{pb:b2:climate-variability:1}
\textbf{1.} $(0.26 + 0.75 + 1.13)/3 = \qty{0.713}{\celsius}$.
\textbf{2.} MJJ $(0.75 + 1.13 + 1.52)/3 = 1.133$; JJA $(1.13 + 1.52 + 1.80)/3 = 1.483$; JAS $(1.52 +
1.80 + 2.06)/3 = 1.793$.
\textbf{3.} ASO $(1.80 + 2.06 + 2.26)/3 = 2.040$; SON $(2.06 + 2.26 + 2.42)/3 = 2.247$; OND $(2.26
+ 2.42 + 2.36)/3 = 2.347$; NDJ $(2.42 + 2.36 + 2.33)/3 = 2.370$.
\textbf{4.} Dibulatkan ke satu desimal, 0.7, 1.1, 1.5, 1.8, 2.0, 2.2, 2.3, dan 2.4: sama dengan tabel.
\textbf{5.} Satu bulan tunggal membawa derau cuaca; ENSO berlangsung berbulan-bulan, dan rata-rata tiga bulan
menghilangkan sebagian besar derau itu sambil mempertahankan sinyal yang lambat.
\textbf{6.} Pada JJA 1997, dengan 1.483, dibulatkan menjadi 1.5; puncaknya 2.4 pada NDJ.
\textbf{7.} Dari ASO 2018 sampai MJJ 2019, dan dari SON 2019 sampai FMA 2020.
\textbf{8.} Sepuluh musim, dan enam.
\textbf{9.} Pada ASO 2020; di dalam tahun 2020 hanya empat musim, ASO sampai NDJ, yang sama dengan atau di bawah $-0.5$, kurang
satu dari aturan itu, dan DJF 2021 melengkapinya.
\textbf{10.} Satu El~Niño, SON 2018 sampai MAM 2019, tujuh musim (AMJ 2019 turun ke 0.37); La~Niña
dimulai pada AMJ 2020 ($-0.55$), empat musim sebelum ONI menempatkannya, pada ASO.
\textbf{11.} Tidak: indeks relatif tidak pernah melebihi 0.29 pada akhir 2019. Seluruh samudra tropis
saat itu lebih hangat daripada periode dasarnya sekitar setengah derajat; wilayah Niño-3.4 hangat secara
mutlak, tetapi hampir tidak lebih hangat daripada daerah tropis di sekitarnya.
\textbf{12.} Sebelas.
\textbf{13.} 8, 7, 10, 4, 5, 6, 12, 5, 9, dan 3 tahun.
\textbf{14.} Yang terpendek 3 tahun (2023 sampai 2026), yang terpanjang 12 (1997 sampai 2009).
\textbf{15.} $69/10 = 6.9$ tahun dengan kedua cara: jumlah selang yang berturutan adalah waktu dari
awal pertama sampai awal terakhir, $2026 - 1957 = 69$ tahun.
\textbf{16.} Peristiwa itu sudah berlangsung lima musim, MAM sampai JAS, sama dengan atau di atas \qty{0.5}{\celsius}, dan
mencapai 2.2, di atas 1.5; awalnya, MAM 2026, adalah musim pertama deretan itu, dan musim-musim berikutnya
hanya dapat memperpanjang deretan itu, bukan menggeser awalnya.
\textbf{17.} Tidak, dan memang tidak perlu: ke-25 episode dari semua kekuatan datang sekitar setiap tiga tahun, seperti kata FAQ itu;
episode yang kuat adalah sebagian darinya, sehingga lebih jarang, dan selangnya berkisar dari 3 sampai 12 tahun.
\textbf{18.} $\tau = 1025 \times 3992 \times 50/16.1 = \qty{1.27e7}{s} = \qty{147}{\day}$.
\textbf{19.} Lapisan campuran: $1025 \times 3992 \times 50 = \qty{2.05e8}{J.m^{-2}.K^{-1}}$; kolom udara:
$(10^5/9.81) \times 1004.5 = 10\,194 \times 1004.5 = \qty{1.02e7}{J.m^{-2}.K^{-1}}$, 20 kali
lebih kecil.
\textbf{20.} $1.024 \times 10^7/(1025 \times 3992) = \qty{2.5}{m}$ air laut.
\textbf{21.} $1.024 \times 10^7/16.1 = \qty{6.36e5}{s} = 7.4$~hari.
\textbf{22.} Anomali udara saja teredam dalam sekitar sepekan, anomali lapisan campuran dalam
sekitar lima bulan, dan lapisan hangat di atas termoklin lebih dalam lagi: hanya samudra yang menyimpan
kalor cukup lama untuk membawa anomali dari satu musim ke musim berikutnya, sedangkan atmosfer menyediakan
angin yang mengopel keduanya.
\textbf{23.} Sejak 1950, episode El~Niño kuat dimulai rata-rata setiap \textbf{6.9 tahun},
sekitar tujuh.
\end{solution}
